\section{準備と問題定義}
\label{sec:preliminaries}

\subsection{効用付きトランザクションデータベース}

\begin{definition}[効用トランザクションデータベース]
\emph{効用トランザクションデータベース} $\mathcal{D} = \{T_0, T_1, \ldots, T_{n-1}\}$ はトランザクションの順序付き列である。各トランザクション$T_t$は時間位置$t$に対応し、アイテム集合$I(T_t) \subseteq \mathcal{I}$を含む。ここで$\mathcal{I}$はアイテムの全体集合である。各アイテム$i \in I(T_t)$は\emph{内部効用}（数量）$q(i, T_t) \in \mathbb{Z}_{>0}$を持つ。各アイテム$i \in \mathcal{I}$は\emph{外部効用}（単位利益）$p(i) \in \mathbb{R}_{>0}$を持つ。
\end{definition}

\begin{definition}[アイテムおよびアイテムセットの効用]
トランザクション$T_t$におけるアイテム$i$の\emph{効用}は$u(i, T_t) = q(i, T_t) \cdot p(i)$である。トランザクション$T_t$におけるアイテムセット$X \subseteq \mathcal{I}$の\emph{効用}は：
\begin{equation}
    u(X, T_t) = \begin{cases}
        \sum_{i \in X} u(i, T_t) & \text{if } X \subseteq I(T_t) \\
        0 & \text{otherwise}
    \end{cases}
\end{equation}
$T_t$の\emph{トランザクション効用}は$\text{TU}(T_t) = \sum_{i \in I(T_t)} u(i, T_t)$である。
\end{definition}

\subsection{スライディングウィンドウと密集区間}

\begin{definition}[スライディングウィンドウ]
ウィンドウサイズ$W \in \mathbb{Z}_{>0}$に対して、位置$l$における\emph{ウィンドウ}はサブデータベース$\mathcal{D}[l, l{+}W] = \{T_t \mid l \leq t \leq l + W\}$である。
\end{definition}

\begin{definition}[ウィンドウ支持度]
位置$l$におけるアイテムセット$X$の\emph{ウィンドウ支持度}は：
\begin{equation}
    \sigma_W(X, l) = |\{T_t \in \mathcal{D}[l, l{+}W] \mid X \subseteq I(T_t)\}|
\end{equation}
\end{definition}

\begin{definition}[密集区間]
アイテムセット$X$は、ウィンドウサイズ$W$および頻度閾値$\theta_f$に対して、すべての$l \in [s, e]$で$\sigma_W(X, l) \geq \theta_f$を満たすとき、\emph{密集区間}$[s, e]$を持つ。
密集区間$[s, e]$が\emph{極大}であるとは、同じ条件を満たす真の拡張$[s', e'] \supsetneq [s, e]$が存在しないことをいう。
\end{definition}

\subsection{ウィンドウ効用とTWU}

\begin{definition}[ウィンドウ効用]
\label{def:wu}
位置$l$におけるアイテムセット$X$の\emph{ウィンドウ効用}は：
\begin{equation}
    \text{WU}(X, l) = \sum_{t=l}^{l+W} u(X, T_t)
\end{equation}
ここで$u(X, T_t)$は定義2より、$X \not\subseteq I(T_t)$の場合は0であることに注意する。したがってWUは、ウィンドウ内で$X$が出現するトランザクションのみから効用を累積する。
\end{definition}

\begin{definition}[ウィンドウトランザクション加重利用度]
\label{def:wtwu}
位置$l$におけるアイテムセット$X$の\emph{ウィンドウTWU}は：
\begin{equation}
    \text{WTWU}(X, l) = \sum_{\substack{t=l \\ X \subseteq I(T_t)}}^{l+W} \text{TU}(T_t)
\end{equation}
直感的には、WTWUはウィンドウ内で$X$を含むトランザクション全体の効用を合算した値であり、WUの上界として機能する。
\end{definition}

\subsection{問題定義}

\begin{definition}[効用密集区間]
\label{def:udi}
アイテムセット$X$は、パラメータ$(W, \theta_f, \theta_u)$に対して、すべての$l \in [s, e]$で以下を\emph{同時に}満たすとき、\emph{効用密集区間}（UDI）$[s, e]$を持つ：
\begin{enumerate}
    \item $\sigma_W(X, l) \geq \theta_f$ \quad （頻度条件）
    \item $\text{WU}(X, l) \geq \theta_u$ \quad （効用条件）
\end{enumerate}
\end{definition}

\begin{definition}[極大効用密集区間]
\label{def:maximal-udi}
効用密集区間$[s, e]$が\emph{極大}であるとは、$[s', e'] \supsetneq [s, e]$なる効用密集区間が存在しないことをいう。すなわち、$s' < s$ならば位置$s'-1$で頻度条件または効用条件が破れ、$e' > e$ならば位置$e'+1$で同様に条件が破れる。
\end{definition}

\textbf{問題定式化。}
効用トランザクションデータベース$\mathcal{D}$、ウィンドウサイズ$W$、頻度閾値$\theta_f$、効用閾値$\theta_u$、最大アイテムセット長$k_{\max}$が与えられたとき、$[s, e]$が$X$の極大効用密集区間（定義~\ref{def:maximal-udi}）であり$|X| \leq k_{\max}$を満たすすべてのペア$(X, [s, e])$を求めよ。
